% ============================================================
% Rancangan Paper Kolaborasi: Edge AI Monitoring Fermentasi Kombucha
% Draft diskusi v0.1 -- Agustus 2026
% Kompilasi: tectonic outline.tex
% ============================================================
\documentclass[conference]{IEEEtran}
\usepackage{amsmath,amssymb,amsfonts}
\usepackage{textcomp}
\usepackage{graphicx}
\usepackage{booktabs}
\usepackage{tabularx}
\usepackage{hyperref}
\usepackage{array}

\hypersetup{
    colorlinks=true,
    linkcolor=black,
    urlcolor=blue,
    citecolor=black
}

\begin{document}

\onecolumn

\title{Rancangan Paper Kolaborasi: Edge AI Monitoring Fermentasi Kombucha (Draft Diskusi v0.1)}
\author{\IEEEauthorblockN{Fatih Jawwad Al Mumtaz (Informatika) \& Amirul Al Aziz (Teknologi Industri Pertanian)}\\
\IEEEauthorblockA{University of Darussalam Gontor, Ponorogo, Indonesia\\
Status: draft diskusi -- bukan draft paper final}
}
\maketitle

\section{Ringkasan Ide}
Membangun sistem monitoring fermentasi kombucha berbasis edge AI \emph{low-cost}: kombinasi sensor murah (pH, Brix, suhu) + model machine learning yang jalan langsung di perangkat (Raspberry Pi 5 / ESP32), untuk memprediksi status dan kematangan fermentasi secara real-time. Sisi Teknologi Industri Pertanian (TIP) membawa domain fermentasi dan protokol data; sisi Informatika membawa pipeline ML dan deployment edge.

\section{Kenapa Proyek Ini Layak (Novelty \& Posisi)}
\begin{itemize}
    \item Literatur ML + kombucha sudah ada tapi semua \textbf{lab-scale}: ANN untuk optimasi fermentasi (2022), prediksi yield SCOBY dari konsentrasi teh/gula/SCOBY/inokulum/pH/suhu/waktu (2022, 2025), review chapter ML di kombucha (2025).
    \item \textbf{Gap:} belum ada yang melakukan monitoring real-time low-cost \emph{on-device} dengan sensor murah. Itu persis niche edge AI yang sudah Fatih bangun (benchmark Burn/TFLite/ONNX di ARM64, protokol pengukuran standar).
    \item Nilai interdisipliner: pangan + AI. Venue IEEE dengan scope CS biasanya terbuka untuk aplikasi interdisipliner selama metodenya solid.
\end{itemize}

\section{Research Questions (RQ)}
\begin{enumerate}
    \item \textbf{RQ1:} Parameter sensor murah mana (pH, Brix, suhu, total asam tertitrasi) yang paling prediktif terhadap status fermentasi?
    \item \textbf{RQ2:} Seberapa akurat model bisa memprediksi kematangan (``siap konsumsi'') dibandingkan metode lab standar (titrasi)?
    \item \textbf{RQ3:} Apakah model tetap akurat setelah dikuantisasi dan di-deploy ke perangkat edge murah (latency/akurasi trade-off)?
    \item \textbf{RQ4 (opsional):} Bisa tidak model membantu optimasi resep batch berikutnya (gula/teh/inokulum)?
\end{enumerate}

\section{Metodologi (Draft -- Butuh Input Sisi TIP)}
\subsection{Protokol Data Collection (pemilik: TIP)}
\begin{itemize}
    \item \textbf{Variabel usulan:} pH, Brix, suhu, total asam tertitrasi (TAT), berat SCOBY, waktu fermentasi.
    \item \textbf{Alat:} pH meter murah, refraktometer, termometer; titrasi manual untuk ground truth. \emph{[TBD: alat yang tersedia di lab kampus, budget]}
    \item \textbf{Skema:} sampling tiap 12--24 jam selama 10--14 hari; minimal 3--5 batch dengan variasi resep (konsentrasi gula/teh/inokulum); ada batch kontrol.
    \item \textbf{Dokumentasi:} log terstruktur (CSV) + foto kondisi batch per sampling.
\end{itemize}

\subsection{Pipeline ML (pemilik: Informatika)}
\begin{itemize}
    \item \textbf{Fitur:} time-series sensor (pH, Brix, suhu, TAT) + metadata batch (resep, suhu ruang).
    \item \textbf{Target:} regresi (trajektori pH/TAT) dan klasifikasi (readiness: siap/tidak, berdasar threshold yang disepakati).
    \item \textbf{Model:} baseline gradient boosting; lalu model sederhana (MLP kecil) yang bisa dikuantisasi.
    \item \textbf{Evaluasi:} RMSE/MAE untuk regresi, accuracy/F1 untuk klasifikasi. \textbf{Penting:} split data per batch (bukan acak) untuk menghindari data leakage.
\end{itemize}

\subsection{Edge Deployment (pemilik: Informatika)}
\begin{itemize}
    \item Device: Raspberry Pi 5 (logging + inferensi) atau ESP32 (jika cukup ringan).
    \item Framework: Burn (Rust) atau TFLite; kuantisasi (int8/fp16).
    \item Evaluasi latency/energi memakai protokol edge benchmark yang sudah ada (1000 warm-up + 1000 run, outlier IQR, bootstrap CI 95\%, seed tetap).
\end{itemize}

\section{Pembagian Kerja (Draft)}
\begin{table}[h]
\centering
\begin{tabular}{@{}lll@{}}
\toprule
\textbf{Aktivitas} & \textbf{Pemilik} & \textbf{Waktu} \\
\midrule
Protokol data \& persiapan lab & TIP & Agu--Sep \\
Batch fermentasi + logging & TIP & Sep--Okt \\
Pra-proses + training ML & Informatika & Okt--Nov \\
Edge deploy + evaluasi & Informatika & Nov \\
Draft \& submission & Bareng (leading: Informatika) & Des--Jan \\
\bottomrule
\end{tabular}
\end{table}

\section{Timeline (Milestone)}
\begin{table}[h]
\centering
\begin{tabular}{@{}ll@{}}
\toprule
\textbf{Periode} & \textbf{Target} \\
\midrule
Agu 2026 & Draft protokol final, diskusi, batch pilot (TIP) \\
Sep--Okt 2026 & Data penuh 3--5 batch (TIP); paralel: MCSoC R2 + Paper A (Informatika) \\
Okt--Nov 2026 & ML training + edge deployment (Informatika) \\
Des 2026 & Draft lengkap siap review bareng \\
Jan--Feb 2027 & Submit ke venue IEEE (cek deadline) \\
\bottomrule
\end{tabular}
\end{table}

\section{Venue Kandidat (Perlu Cek Deadline)}
\begin{itemize}
    \item IEEE Indonesia: ISRITI, ICACSIS, ICoDSE, EECSI (umumnya IEEE Xplore + Scopus).
    \item Opsi internasional: IEEE venue food-tech/agro-AI jika ada yang sesuai (perlu riset).
    \item Kriteria: IEEE Xplore + Scopus/Ei, deadline fleksibel, \emph{turnaround} cepat (konsisten dengan preferensi venue yang sudah ada).
\end{itemize}

\section{Keputusan yang Harus Didiskusikan (Pertanyaan Terbuka)}
\begin{enumerate}
    \item \textbf{First authorship \& urutan author.} Usulan: yang menulis inti paper dan analisis jadi first author. Perlu kesepakatan eksplisit sebelum kerja mulai.
    \item \textbf{Protokol data:} alat apa yang benar-benar tersedia di lab TIP? Berapa budget sensor yang realistis?
    \item \textbf{Standar ``siap konsumsi'':} ditentukan oleh apa (pH? keasaman? sensorik?)? Butuh literatur food science dari sisi TIP.
    \item \textbf{Plan B:} kalau akurasi model kurang, scope turun ke \emph{feasibility study} atau ganti target (misal prediksi yield SCOBY)?
    \item \textbf{Halal angle:} masuk sekarang (butuh spektroskopi, mahal) atau dijadikan fase 2?
\end{enumerate}

\section{Langkah Berikutnya (Minggu Ini)}
\begin{enumerate}
    \item Diskusi protokol data (TIP): inventaris alat lab + variabel yang feasible.
    \item Mulai batch pilot sesegera mungkin (TIP) -- data adalah long pole proyek ini.
    \item Cek venue \& deadline (Informatika).
    \item Review literatur singkat masing-masing sisi (2--3 paper per sisi).
\end{enumerate}

\end{document}
